\clearpage\fancyhead[L]{\small JUNO.}\hypertarget{native-377-lower}{}\label{native-377-lower}
\begin{center}\textbf{JUNO.}\par\vspace{4pt}\hrule\vspace{4pt}\fontsize{10}{12}\selectfont Göttingische gelehrte Anzeigen. Stück 204. Seite 2033\;..\;2034. 23. Dezember 1813.\par\vspace{4pt}\hrule\vspace{4pt}\end{center}


\par\indent Die Jahreszeit, in der \emph{Juno} dieses Jahr in der \emph{Georgsharfe} in Opposition zur Sonne kam, ist hier gewöhnlich ungünstig für astronomische Beobachtungen. Professor \textsc{Gauss} konnte um diese Zeit nur eine einzige Beobachtung dieses Planeten erhalten:
\begin{center}\fontsize{10}{12}\selectfont\setlength{\tabcolsep}{4pt}\renewcommand{\arraystretch}{1.15}\begin{tabular}{l|c|c|c}
1813 & Mittlere Zeit & Rektaszension & Deklination\\\hline
Nov. 19 & $9^{\mathrm u}\;46^{\mathrm m}\;46''$ & $60^\circ\;30'\;35''0$ & $3^\circ\;5'\;9''3\ \text{Südl.}$\\
\end{tabular}\end{center}
\par\indent Obwohl der Planet diesmal zugleich sehr nahe an seinem Perihel und damit hinsichtlich seiner Beleuchtung in der günstigsten Lage war, erreichte er doch kaum die 8. Größe.

\Needspace{14\baselineskip}\par\indent Da auswärtige Beobachtungen ganz fehlen und derzeit von entfernten Sternwarten keine zu erwarten sind, selbst wenn der Planet noch irgendwo anders beobachtet worden sein sollte, musste sich Professor \textsc{Gauss} entschließen, die Berechnung der Opposition auf die obigen einzelnen Beobachtungen zu stützen. Das Ergebnis lautet:
\begin{center}\leavevmode\SourceLeaf{378}\emph{Achte Opposition der Juno.}\end{center}
\begin{center}1813 November 19. $18^{\mathrm u}\;11^{\mathrm m}\;5^{\mathrm s}$ Mittlere Zeit in Göttingen.\end{center}
\begin{center}\fontsize{10}{12}\selectfont\setlength{\tabcolsep}{2pt}\renewcommand{\arraystretch}{1.15}\begin{tabular}{lcr}
Wahre Länge & $.\;.\;.\;.\;.\;.\;.\;.\;.\;.\;.\;.\;.$ & $57^\circ\;33'\;58''0$\\
Geozentrische Breite & $.\;.\;.\;.\;.\;.\;.\;.\;.\;.\;.$ & $23\;18\;46.8\ \text{Südl.}$\\
\end{tabular}\end{center}
\Needspace{20\baselineskip}\par\indent Für die nächste Sichtbarkeitsperiode des Planeten soll eine Ephemeride berechnet und an anderer Stelle mitgeteilt werden. Dafür leitete Professor \textsc{Gauss} aus der fünften, sechsten, siebten und achten Opposition neue elliptische Elemente ab.
\begin{center}\fontsize{9}{11}\selectfont\setlength{\tabcolsep}{2pt}\renewcommand{\arraystretch}{1.15}\begin{tabular}{lcr}
\shortstack[l]{Epoche der mittleren Länge 1810\\im Meridian von Göttingen} &  & $95^\circ\;29'\;56''$\\
Tägliche tropische mittlere Bewegung & $.\;.\;.\;.\;.\;.\;.$ & $812''709$\\
Tropische Umlaufszeit & $.\;.\;.\;.\;.\;.\;.\;.\;.\;.\;.\;.\;.$ & $1594\ \text{Tage}\ 16\ \text{Stunden}$\\
Länge des Perihels (1810) & $.\;.\;.\;.\;.\;.\;.\;.\;.\;.\;.$ & $53^\circ\;6'\;40''$\\
Länge des aufsteigenden Knotens (1810) & $.\;.\;.\;.\;.\;.\;.$ & $171\;7\;13$\\
Exzentrizitätswinkel & $.\;.\;.\;.\;.\;.\;.\;.\;.\;.\;.\;.\;.\;.$ & $14\;43\;8$\\
Neigung der Bahn & $.\;.\;.\;.\;.\;.\;.\;.\;.\;.\;.\;.\;.\;.$ & $13\;4\;12$\\
Log. der großen Halbachse & $.\;.\;.\;.\;.\;.\;.\;.$ & $0.4267670$\\
\end{tabular}\end{center}
\par\indent Der Vergleich dieser Elemente mit denen, die auf den vier \emph{ersten} Oppositionen beruhten (siehe Stück 136 dieser Anzeigen von 1808 [S. 306 dieses Bandes]), zeigt sehr deutlich, wie stark auch bei diesem Planeten die äußeren Störungen wirken. Besonders auffällig ist der Unterschied in der mittleren Bewegung. Aus den früheren Oppositionen ergab sich eine um 3 Tage und 4 Stunden kürzere Umlaufszeit. Ohne die Störungen hätte dieser Unterschied keine halbe Stunde betragen können.
